\DocumentMetadata{
  pdfversion=2.0,
  pdfstandard=ua-2,
  lang=en-US,
  tagging=on,
  tagging-setup={math/setup=mathml-SE,tabsorder=structure}
}
\documentclass[11pt,letterpaper]{article}
\usepackage[margin=0.68in,includefoot]{geometry}
\usepackage{newcomputermodern}
\usepackage{amsmath,amscd,mathtools}
\usepackage{tikz,adjustbox}
\providecommand{\square}{\mdlgwhtsquare}
\usepackage[hidelinks]{hyperref}
\hypersetup{
  pdftitle={Numerical Analysis PhD Exam, May 2001},
  pdfauthor={University of Florida Department of Mathematics},
  pdfsubject={Graduate mathematics examination},
  pdfdisplaydoctitle=true,
  bookmarksopen=true
}
\setcounter{secnumdepth}{0}
\setlength{\parindent}{0pt}
\setlength{\parskip}{0.28em}
\setlength{\emergencystretch}{3em}
\setlength{\leftmargini}{2.15em}
\setlength{\leftmarginii}{2.65em}
\setlength{\leftmarginiii}{3em}
\renewcommand{\labelenumi}{\textbf{\theenumi.}}
\pagestyle{empty}
\raggedbottom
\allowdisplaybreaks
\newenvironment{examproblems}[1]{%
  \begin{enumerate}%
  \setcounter{enumi}{#1}%
  \setlength{\topsep}{0.35em}%
  \setlength{\partopsep}{0pt}%
  \setlength{\itemsep}{0.48em}%
  \setlength{\parsep}{0.18em}%
}{\end{enumerate}}
\newenvironment{examsubparts}{%
  \begin{enumerate}%
  \setlength{\topsep}{0.18em}%
  \setlength{\partopsep}{0pt}%
  \setlength{\itemsep}{0.16em}%
  \setlength{\parsep}{0.1em}%
}{\end{enumerate}}
\newenvironment{examinstructions}{%
  \par\begingroup\itshape
}{%
  \par\endgroup\vspace{0.18em}
}
\makeatletter
\renewcommand\section{\@startsection{section}{1}{0pt}%
  {-1.25ex plus -.25ex minus -.1ex}%
  {0.55ex plus .1ex}%
  {\normalfont\large\bfseries}}
\renewcommand{\@maketitle}{%
  \newpage\null\vskip -1em%
  {\footnotesize\bfseries
    \href{https://www.ufl.edu/}{University of Florida}, \href{https://math.ufl.edu/}{Department of Mathematics}\par}%
  \vskip -0.5em%
  \tagstructbegin{tag=Title}%
    {\LARGE\bfseries\href{https://gma.math.ufl.edu/exams/numerical-analysis-phd/}{Numerical Analysis PhD Exam}, May 2001\par}%
  \tagstructend%
  \vskip 0.25em%
  \tagmcbegin{artifact}%
    \hrule height 1pt%
  \tagmcend%
  \vskip 1em%
}
\makeatother
\title{Numerical Analysis PhD Exam, May 2001}
\author{}
\date{}
\begin{document}
\maketitle
\thispagestyle{empty}
\begin{examinstructions}
Do any 8 of the following 10 problems.
\end{examinstructions}
\section{Part 1: Numerical Linear Algebra}
\begin{examproblems}{0}
\item
This problem has five parts.
\begin{examsubparts}
\item[\textnormal{(a)}]
Under what conditions does a real~\(n\) by~\(n\) matrix~\(A\) have a Schur decomposition \(A=U^{\mathrm T}TU\), where~\(U\) and~\(T\) are real~\(n\) by~\(n\) matrices with~\(U\) orthogonal and~\(T\) upper triangular?
\item[\textnormal{(b)}]
Under what conditions on~\(A\) is~\(T\) diagonal?
\item[\textnormal{(c)}]
Does the matrix 
\[
A=\begin{pmatrix}1&0&1\\1&-1&0\\0&1&1\end{pmatrix}
\]
 have an orthogonal set of eigenvectors?
\item[\textnormal{(d)}]
Give a careful statement of the singular value decomposition of a real~\(m\) by~\(n\) matrix.
\item[\textnormal{(e)}]
Under what conditions does the singular value decomposition exist?
\end{examsubparts}
\item
Consider the linear system \(Ax=b\).
\begin{examsubparts}
\item[\textnormal{(a)}]
If the rows of~\(A\) are linearly independent, derive a formula for the solution of minimal 2-norm.
\item[\textnormal{(b)}]
If the columns of~\(A\) are linearly independent, derive a formula for the~\(x\) that minimizes the 2-norm of the residual \(r=b-Ax\).
\item[\textnormal{(c)}]
Use the singular value decomposition of~\(A\) to give a formula for the~\(x\) that minimizes the 2-norm of the residual, and among all the~\(x\)’s which minimize the 2-norm, it has minimal norm.
\end{examsubparts}
\item
This problem has two parts.
\begin{examsubparts}
\item[\textnormal{(a)}]
Given a vector \(x\in\mathbb R^n\) and a natural number \(k<n\), give a formula for a unit vector~\(w\) for which the vector \(y=(I-2ww^{\mathrm T})x\) satisfies the following conditions: \(y_i=x_i\) for \(i<k\) and \(y_i=0\) for \(i>k\).
\item[\textnormal{(b)}]
Using these Householder transformations, write a pseudo code (or a matlab code) for reducing a real symmetric matrix to tridiagonal form using Householder similarity transformations.
\end{examsubparts}
\item
This problem has three parts.
\begin{examsubparts}
\item[\textnormal{(a)}]
State and prove the Gershgorin Circle Theorem.
\item[\textnormal{(b)}]
Let 
\[
A=\begin{pmatrix}5&1&3\\2&4&1\\3&-1&4\end{pmatrix}.
\]
 What can you say about the location of the eigenvalues of~\(A\)?
\item[\textnormal{(c)}]
Use the Gershgorin Circle Theorem to estimate the size of the largest root of the polynomial \(p(x)=x^n+a_{n-1}x^{n-1}+a_{n-2}x^{n-2}+\cdots+a_1x+a_0\).
\end{examsubparts}
\item
This problem has two parts.
\begin{examsubparts}
\item[\textnormal{(a)}]
Find the ellipse or hyperbola, of the form \(ax^2+by^2=1\), that best fits~\(n\) data points \((x_1,y_1),(x_2,y_2),\ldots,(x_n,y_n)\) in the plane by computing the least squares solution to the~\(n\) linear equations gotten by substituting each of the data points into the quadratic equation.
\item[\textnormal{(b)}]
Test your method with the dataset \((1,1),(2,4),(3,9),(4,16)\).
\end{examsubparts}
\end{examproblems}
\section{Part II: Numerical Analysis}
\begin{examproblems}{5}
\item
This problem has three parts.
\begin{examsubparts}
\item[\textnormal{(a)}]
State Newton’s method (the algorithm) for solving \(f(x)=0\) where \(f:\mathbb R\to\mathbb R\).
\item[\textnormal{(b)}]
Assuming~\(f\) is smooth and we have an initial guess sufficiently close to a root~\(p\), state sufficient condition(s) for the method to converge quadratically.
\item[\textnormal{(c)}]
Show that Newton’s method applied to \(x^5=2\) converges quadratically to the root \(2^{1/5}\) from any starting guess \(x>0\).
\end{examsubparts}
\item
This problem has three parts.
\begin{examsubparts}
\item[\textnormal{(a)}]
Given a function~\(f\) with \(n+1\) continuous derivatives on the interval \(I=[-1,1]\), show that, given \(n+1\) points \(\{x_0,x_1,\ldots,x_n\}\) in~\(I\), there exists a unique polynomial \(P_n(x)\) of degree \(\le n\) such that 
\[
P_n(x_i)=f(x_i)\qquad\text{for }i=0,\ldots,n.
\]
\item[\textnormal{(b)}]
Prove that, given \(t\in I\), there exists \(\eta\in I\) such that 
\[
f(t)-P_n(t)=\frac{(t-x_0)\cdots(t-x_n)}{(n+1)!}f^{(n+1)}(\eta).
\]
\item[\textnormal{(c)}]
From the above we get that (all norms are on the interval~\(I\)) 
\[
\lVert f-P_n\rVert_\infty\le \max_{t\in I}\left|(t-x_0)(t-x_1)\cdots(t-x_n)\right|\frac{\lVert f^{(n+1)}\rVert_\infty}{(n+1)!}.
\]
 What choice for the points \(x_0,\ldots,x_n\) minimizes the right hand side of this error bound?
\end{examsubparts}
\item
This problem has two parts.
\begin{examsubparts}
\item[\textnormal{(a)}]
Find a function \(q(x)\) which has a continuous first derivative for each \(x\in(-\infty,\infty)\) with the properties: 
\[
q(x)=\begin{cases}
1,&\text{if }x=0,\\
0,&\text{for all }|x|\ge 2,\\
a_0+a_1x+a_2x^2,&\text{for all }x\in[-2,-1],\\
b_0+b_1x+b_2x^2,&\text{for all }x\in[-1,1],\\
c_0+c_1x+c_2x^2,&\text{for all }x\in[1,2].
\end{cases}
\]
\item[\textnormal{(b)}]
Use the function \(q(x)\) defined in part a to interpolate a set of equally spaced points \((x_0,y_0),(x_1,y_1),\ldots,(x_n,y_n)\), where \(h=x_{k+1}-x_k\) for all~\(k\). In particular, describe how to find constants \(c_k\) so that the function 
\[
Q(x)=\sum_{k=0}^n c_k q\!\left(\frac{x-x_k}{h}\right)
\]
 has the property that \(Q(x_k)=y_k\) for all~\(k\).
\end{examsubparts}
\item
Let \(P_2(x)\) be a quadratic polynomial interpolating \(g(x)\) at \(x=0,h,2h\).
\begin{examsubparts}
\item[\textnormal{(a)}]
Use this to derive a numerical integration formula \(I_h\) for 
\[
I=\int_0^{3h}g(x)\,dx.
\]
\item[\textnormal{(b)}]
Use a Taylor series expansion of \(g(x)\) to show 
\[
I-I_h=\frac{3}{8}h^4g^{(3)}(0)+O(h^5).
\]
\end{examsubparts}
\item
Triple Recursion Formula. If \(\{\phi_n(x)\}\) is an orthogonal family of polynomials on \([a,b]\), with respect to the weight function \(w(x)\ge 0\), and \(n\ge 1\), then show that 
\[
\phi_{n+1}(x)=(a_nx+b_n)\phi_n(x)-c_n\phi_{n-1}(x),
\]
 for some constants \(a_n,b_n\), and \(c_n\). (Hint: Let \(g(x)=\phi_{n+1}(x)-a_nx\phi_n(x)\), where \(a_n\) is a constant chosen so that \(g(x)\) has degree~\(n\).)
\end{examproblems}
\end{document}
